\graphicspath{{3.nal_slike/}}

\section{Schr{\"o}dingerjeva enačba}

Za vodikov atom se Shr{\"o}dingerjeva enačba za radialni del v brezdimenzijski obliki glasi
\begin{equation}
    \left( -\frac{d^2}{dx^2} - \frac{2}{x} + \frac{l\,(l+1)}{x^2} - E\right)\; R(x) = 0, 
\end{equation}
pri čemer se robni pogoji glasijo: $R(x=0)=R(x=\infty)=0$, strelski parameter pa je $E$ in hkrati predstavlja lastno energijo. Za uporabo klasičnih integratorjev, kot je Runge-Kutta, je potrebno zgornjo DE 2. reda prevesti v sistem dveh DE 1. reda
\begin{align}
    \frac{dR}{dx} &= z\\\nonumber
    \frac{dz}{dx} &= \left(-\frac{2}{x} + \frac{l\,(l+1)}{x^2} - E\right)\; R(x),
\end{align}
ki ju nato integriramo v radialni smeri z začetnima pogojema: robnim pogojem kot začetno vrednostjo in začetnim naklonom, ki pa ne vpliva na obliko rešitve, nego le na skaliranje, zato vsako rešitev normiram, da velja
\begin{equation}
    \int_{0}^{\infty} R^2(x)\,dx = 1.
\end{equation}
Še boljša pa je metoda Numerova, kjer enačbo oblike
\begin{equation}
    \left( \frac{d^2}{dx^2} + k^2(x) \right)\; y(x) = 0,
\end{equation}
s funkcijo $k$ za ta problem
\begin{equation}
    k(x) = \frac{2}{x} - \frac{l\,(l+1)}{x^2} + E,
\end{equation}
zapišemo v diskretni obliki s $h$ kot časovnim korakom
\begin{equation}
    y_{i+1} = \frac{2\,\left(1-\frac{5\,h^2}{12}\,k_i^2\right) \,y_i - \left(1-\frac{h^2}{12}\,k_{i-1}^2 \right)\,y_{i-1}}{1-\frac{h^2}{12}\,k_{i+1}^2},
\end{equation}
kjer sem uporabil začetne pogoje, enake kot prej, vendar Numerov potrebuje dve začetni vrednosti, zato sem drugo pridobil z linearno aproksimacijo iz začetnih pogojev
\begin{equation}
    y_1 = y_0 + h\,z_0,
\end{equation}
kjer je $z_0$ začetni naklon, kateri zaradi normiranja rešitve nima fizikalnega pomena.

Pri integraciji naprej sem začel pri $x=0$ z začetnimi pogoji in integriral do $x_{max}$, pri metodi nazaj pa ravno obratno. Obojestransko sem integriral do $x_{max}/2$, nakar sem rešitev skaliral tako, da je bil naklon z obeh strani enak. Potrebno je bilo še prilagoditi strelski parameter $E$ tako, da bo izpolnjen robni pogoj na drugi strani oziroma, da se integracija iz obeh strani sreča v eni točki. Za iskanje ničle te napake lege sem uporabil sekantno metodo z začetnima približkoma, blizu analitični vrednosti za $E$, $x_0= 1/n^2$ in $x_1=1.001\,x_0$.

Zaradi narave lastnih funkcij $R(x)$ je bilo potrebno prilagajati tudi $x_{max}$, za kar sem uporabil preprosto aproksimacijo, ki deluje za $n<50$
\begin{equation}
    x_{max} \approx 10\,n^{1.65}.
\end{equation}

Spodaj so prikazane nekatere lastne funkcije $R(x)$, dobljene z metodo Numerov.

\fig{H}{1}{lastne_funk_n}{Lastne funkcije za $l=0$}{}
\fig{H}{1}{lastne_funk_l}{Lastne funkcije z $l>0$}{}

Za zgornje primere sem preveril vrednosti strelskega parametra $E$, ki se precej dobro sklada z analitično zvezo $E=1/n^2$

\fig{H}{0.75}{energije_n}{Izračunane lastne energije}{}

Sledi prikaz absolutne vrednosti napake lege, kjer sem iskal ničlo. Opazim, da se za manjše $x$ vrednost $E$ lepo sklada z analitično rešitvijo, pri višjih pa je verjetno težava z $x_{max}$. Zgleda, da se najbolje se izkaže še obojestranska integracija, sledi metoda nazaj in najslabša naprej.

\fig{H}{0.75}{grid_search}{\textit{Grid search} za strelski parameter energije $E$. Navpično so označene analitične vrednosti.}{}

\begin{table}[H]
    \begin{tabular}{c|c c c c c}
        metoda & RK naprej & RK nazaj & RK obojestransko & Numerov naprej & Numerov nazaj \\\hline
        čas [s] & 0,9 & 3,1 & 1,0 & 0,2 & 0,4
    \end{tabular}
    \caption{Časi izračunov za \textit{Grid search} -- 1000 izračunov za $E \in (10^{0.05}, 10^{-2.5})$}
\end{table}

Iz sledečih slik lahko razberemo, v kolikšni meri se numerična rešitev $R_{nl}$ sklada z analitično $\mathcal{R}_{nl}$ za različne metode. V neskončnosti zaradi numeričnih težav predpišem vrednost blizu nič ($10^{-4}$), prav tako tudi ne integriram točno do $x=0$, ampak do $x=10^{-9}$. Ker se v primeru sodega števila ničel v $x=\infty$ ničli približa od spodaj, moram to upoštevati tam pri robnem pogoju.

Opazim, da se z zmanjšanejem $\Delta x$, napak manjša, razen za integrator RK45 iz \\\texttt{scipy.integrate.solve\_ivp}, kjer je verjetno težava, in hkrati očitna pohitritev v adaptivnem koraku. Metoda Numerov pričakovano dosega najnižje napake, obojestranska pa ima 

Nekako je metoda nazaj 2x počasnejša, vendar je bolj stabilna, saj se strelja proti vrednosti $R(x=0)=0$, dočim pri metodi naprej ta cilj ni nikoli dosežen, saj se zaradi numerične napake v neskončnosti oddaljujemo od željenega robnega pogoja.

To se sicer lahko kompenzira z dovolj majhnim $x_{max}$, vendar je že zagotavljanje dovolj velikega, da zajame vse ničle, povzročalo težavo, ali ustavitvijo integracije pri določeni bližini x-osi, kar je lahko težavno zaradi večjega števila ničel.

\begin{figure}[H]
    \centering
    \subfig{0.48}{napaka_dx_n1_l0_naprej-RK}{lastna koda RK4 naprej}{}
    \subfig{0.48}{napaka_dx_n1_l0_naprej}{RK45 naprej}{}
    \subfig{0.48}{napaka_dx_n1_l0_nazaj}{RK45 nazaj}{}
    \subfig{0.48}{napaka_dx_n1_l0_obojestransko}{RK45 obojestransko}{}
    \subfig{0.48}{napaka_dx_n1_l0_numerov}{Numerov naprej}{}
    \subfig{0.48}{napaka_dx_n1_l0_numerov nazaj}{Numerov nazaj}{}
    \caption{Vpliv $\Delta x$ na napako rešitve $R_{10}(x)$}
\end{figure}

\begin{figure}[H]
    \centering
    \subfig{0.48}{napaka_dx_n2_l1_numerov nazaj}{Numerov nazaj}{}
    \subfig{0.48}{napaka_dx_n2_l1_obojestransko}{obojestransko}{}
    \caption{Vpliv $\Delta x$ na napako rešitve $R_{21}(x)$}
\end{figure}

\begin{figure}[H]
    \centering
    \subfig{0.48}{napaka_dx_n3_l0_numerov nazaj}{Numerov nazaj}{}
    \subfig{0.48}{napaka_dx_n3_l0_obojestransko}{obojestransko}{}
    \caption{Vpliv $\Delta x$ na napako rešitve $R_{30}(x)$}
\end{figure}

\begin{figure}[H]
    \centering
    \subfig{0.48}{napaka_dx_n3_l1_numerov nazaj}{Numerov nazaj}{}
    \subfig{0.48}{napaka_dx_n3_l1_obojestransko}{obojestransko}{}
    \caption{Vpliv $\Delta x$ na napako rešitve $R_{31}(x)$}
\end{figure}

\begin{figure}[H]
    \centering
    \subfig{0.48}{napaka_dx_n3_l2_numerov nazaj}{Numerov nazaj}{}
    \subfig{0.48}{napaka_dx_n3_l2_obojestransko}{obojestransko}{}
    \caption{Vpliv $\Delta x$ na napako rešitve $R_{32}(x)$}
\end{figure}

Peizkusil sem še, kako pravilne so rešitve. To sem izvedel tako, da sem preštel ničle fukncije in ničle odvodov, kjer mora obojih posebej biti natanko $n-l$. Očitno se z zmanjšanjem koraka $\Delta x$ stanje izboljša, preseneča pa me metoda Numerov nazaj, kjer je to izboljšanje zelo majhno! Morebiti je temu kriva linearna aproksimacija začetne točke, vendar tudi višji red ni bistveno izboljšal zadeve.

\begin{figure}[H]
    \centering
    \subfig{0.48}{napake_dx0.1_naprej-RK}{lastna koda RK4 naprej}{}
    \subfig{0.48}{napake_dx0.1_naprej}{RK45 naprej}{}
    \subfig{0.48}{napake_dx0.1_nazaj}{RK45 nazaj}{}
    \subfig{0.48}{napake_dx0.1_obojestransko}{RK45 obojestransko}{}
    \subfig{0.48}{napake_dx0.1_numerov}{Numerov naprej}{}
    \subfig{0.48}{napake_dx0.1_numerov nazaj}{Numerov nazaj}{}
    \caption{Pravilnost števila ničel in ničel odvodov pri $\Delta x = 0.1$}
\end{figure}

\begin{figure}[H]
    \centering
    \subfig{0.48}{napake_dx0.01_naprej-RK}{lastna koda RK4 naprej}{}
    \subfig{0.48}{napake_dx0.01_naprej}{RK45 naprej}{}
    \subfig{0.48}{napake_dx0.01_nazaj}{RK45 nazaj}{}
    \subfig{0.48}{napake_dx0.01_obojestransko}{RK45 obojestransko}{}
    \subfig{0.48}{napake_dx0.01_numerov}{Numerov naprej}{}
    \subfig{0.48}{napake_dx0.01_numerov nazaj}{Numerov nazaj}{}
    \caption{Pravilnost števila ničel in ničel odvodov pri $\Delta x = 0.01$}
\end{figure}

Glede porabe časa sem opazil, da je metoda Numerov bistveno hitrejša. Opazna pa je razlika med lastno in scipy Runge-Kutto, sploh pri manjših $\Delta x$.

\fig{H}{0.75}{cas}{Povprečna časovna zahtevnost metod}{}

\section{Helmholtzova enačba}

Rešujem Helmholtzovo enačbo za propagacijo monokromatske svetlobe v radialni smeri snovi, kot je optično vlakno,
\begin{equation}
    \left( \frac{d^2}{dx^2} + \frac{1}{4\,x^2} + n^2(x)\,k^2 - \lambda^2 \right) \, R(x) = 0,
\end{equation}
kjer je $k$ valovno število in hkrati prosti parameter simulacije, $n(x)$ lomni količnik, $\lambda$ lastna vrednost in $R(x)$ lastna funkcija.

V navodilih je podan potek lomnega količnika v radialni smeri optičnega vlakna
\begin{equation}
    n(x) =
    \begin{cases}
        N_j-\frac{x^2}{2} & x < 1\\
        1 & x \geq 1
    \end{cases},
\end{equation}
kjer je $N_j=2$ lomni količnik jedra. Sam pa definiram še večplastno (dvoplastno jedro) vlakno
\begin{equation}
    n(x) =
    \begin{cases}
        N_j-\frac{1}{2}\,x^2 & x < 1\\
        1.25-0.2\,x+0.2 & 1 \leq x < 2\\
        1 & x \geq 2
    \end{cases}
\end{equation}
\begin{figure}[H]
    \centering
    \subfig{0.48}{lomni}{lomni količni iz navodil}{}
    \subfig{0.48}{lomni2}{nov lomi količnik}{}
\end{figure}

Enačbo rešujem na enak način kot vprejšnjem poglavju z metodama Numerov in Runge-Kutta, tudi robni pogoji so (za fizikalno smiselne rešitve) isti. $x_{max}$ pa je drugačen
\begin{equation}
    x_{max}=\frac{20}{k},
\end{equation}
saj ga je za večje $k$ potrebno zmanjševati.

Izvedem iskanje ničle na enak način kot poprej, vendar kmalu ugotovim, da je v tem primeru to zahtevnejše delo, saj so ničle že iz spodnjih grafov komaj opazne. Prav tako me preseneča, da z različnimi metodami dobim različne lege ničel, torej različne vrednosti parametrov $\lambda$! To je bolje opazno v spodnji tabeli.

Ničle napake robnega pogoja sem iskal v dveh korakih. V prvem sem interval, kjer sem opazil, da se nahajajo vse ničle, $\lambda \in (k,2\,k)$, razdelil na 10 podintervalov, na katerih sem opazoval spremembo predznaka. Če je sprememba bila, sem v drugem koraku ničlo na podintervalu natančneje določil z metodo \textit{brent} iz \texttt{scipy.optimize.root\_scalar}. Tako lahko določim max 10 ničel, saj ako sta dve na enem podintervalu algoritem ne zazna nobene!

\begin{figure}[H]
    \centering
    \subfig{0.48}{grid_search_k1}{$k=1$}{}
    \subfig{0.48}{grid_search_k2}{$k=2$}{}
    \subfig{0.48}{grid_search_k5}{$k=5$}{}
    \subfig{0.48}{grid_search_k10}{$k=10$}{}
    \caption{Iskanje ničle na intervalu $\lambda \in (k, 2\cdot k)$}
\end{figure}

\begin{table}[H]
    \begin{tabular}{c|c c c c c c}
        metoda & \shortstack{RK4 lastna\\naprej} & \shortstack{RK4 lastna\\nazaj} & \shortstack{RK45 scipy\\naprej} & \shortstack{RK45 scipy\\nazaj} & \shortstack{Numerov\\naprej} & \shortstack{Numerov\\nazaj} \\ \hline
        $\lambda$ & 1,068 & 1,144 & 1,160 & 1,224 & 1,121 & 1,121\\
    \end{tabular}
    \caption{Rezultati različnih metod pri $k=1$}
\end{table}

Sledi prikaz lastnih oblik za nekaj primerov in disperzijska relacija dva različna lomna količnika jedra.

\begin{figure}[H]
    \centering
    \subfig{0.48}{modes_k1}{$k=1$}{}
    \subfig{0.48}{modes_k2}{$k=2$}{}
    \subfig{0.48}{modes_k3}{$k=3$}{}
    \subfig{0.48}{modes_k5}{$k=5$}{}
    \subfig{0.48}{modes_k7}{$k=7$}{}
    \subfig{0.48}{modes_k10}{$k=10$}{}
    \caption{Lastne funkcije in pripadajoče vrednosti $\lambda$}
\end{figure}

Opazim, da $N_j$ spreminja naklon disperzijske relacije in zniža vrednost valovnega števila $k$, do katerega je vlakno še ``single mode fiber''. Pri $N_j$ je to do $k=2.865$. Prav tako se zniža na $k=2.302$ pri večplastnem vlaknu.

\fig{H}{0.8}{disperzijska}{Disperzijska relacija pri $N_j=2$}{}
\fig{H}{0.8}{disperzijska_2}{Disperzijska relacija pri $N_j=3$}{}
\fig{H}{0.8}{disperzijska_lomni2}{Disperzijska relacija za lomni količnik po večplastnem vlaknu}{}

Na spodnjem grafu je bolje viden vpliv lomnega količnika jedra $N_j$ na mejni $k$ za ``single mode fibre''. Hkrati sem računal lego presečišča premic, vendar, ker je sistem predoločen, verjetno to v nekaterih primerih povzroča težavo, zato iz tega nekakšne korelacije ni mogoče razbrati. Za rešitev sistema sem uporabil rešitev po metodi najmanjših kvadratov po \texttt{numpy.linalg.lstsq}. Prav tako postane težava, ko sem želel na eno točko ``fitati'' premico!

\fig{H}{0.8}{k_single_mode}{Sprememba meje za ``single mode fibre'' v odvisnosti od lomnega količnika jedra}{}
\fig{H}{0.8}{presecisce}{Presečišče premic disperzijske relacije}{}

\fig{H}{0.8}{cas_helmholtz}{Časovna zahtevnost}{}
Časovna zahtevnost me v primerjavi s prejšnjo preseneča, saj je očitno večja. Domnevam, je pri $k=1$ vrh zaradi prvega klica vseh funkcij, saj nato opazim trend naraščanja. 